% Authored directly from the companion Markdown note; no generated prose.
% Compile from the repository root:
% latexmk -pdf -outdir=/tmp/gamma-review-latex docs/gamma_optimization_paper_note.tex
\documentclass[11pt]{article}
\usepackage[T1]{fontenc}
\usepackage{amsmath,amsthm}
\usepackage{newtxtext,newtxmath}
\usepackage[letterpaper,margin=0.75in]{geometry}
\usepackage{graphicx,microtype,xcolor}
\usepackage{caption,enumitem,fancyhdr,needspace}
\usepackage[colorlinks=true,linkcolor=blue!45!black,urlcolor=blue!55!black]{hyperref}
\definecolor{ink}{HTML}{17334A}
\definecolor{accent}{HTML}{C65525}
\hypersetup{pdftitle={Gamma controls the scales available to readout gradient descent},
  pdfsubject={Mechanism theorem, five-target evidence, and complete proof}}
\graphicspath{{results/checkpoint_D_optimizers/expD36_frozen_gamma_probe/full_sweep/refinements/gamma_factorized_kernel/}}
\newcommand{\repo}{https://github.com/samlayton99/precision-mlps/blob/experiment/frozen-gamma-probe}
\newcommand{\evidence}{\repo/results/checkpoint_D_optimizers/expD36_frozen_gamma_probe/full_sweep/refinements/gamma_factorized_kernel}
\DeclareMathOperator{\sech}{sech}
\DeclareMathOperator{\sign}{sign}
\DeclareMathOperator{\Res}{Res}
\theoremstyle{definition}
\newtheorem{theorem}{Theorem}
\newtheorem{corollary}{Corollary}
\setlength{\parindent}{0pt}
\setlength{\parskip}{6pt plus 1pt minus 1pt}
\setlength{\emergencystretch}{2em}
\setlength{\headheight}{14pt}
\captionsetup{font=small,labelfont=bf,justification=raggedright,singlelinecheck=false}
\setlist[itemize]{leftmargin=1.3em,itemsep=3pt,topsep=3pt}
\pagestyle{fancy}
\fancyhf{}
\fancyhead[L]{\small\color{ink}Gamma and readout optimization}
\fancyhead[R]{\small Collaborator review}
\fancyfoot[C]{\small\thepage}
\renewcommand{\headrulewidth}{0.3pt}

\begin{document}
\thispagestyle{plain}
{\color{ink}\fontsize{23}{27}\selectfont\bfseries
Gamma controls the scales available\newline
to readout gradient descent\par}
\vspace{4pt}
{\small\color{ink}Research note \quad Mechanism, measured delays, and self-contained proof}
\vspace{6pt}

\textbf{Gamma changes how quickly a readout learns an attainable target.}
For fixed neuron centers, a common slope applies an explicit frequency filter
to the features and hence to the kernel governing readout updates. Small gamma
strongly attenuates fine spatial variation. Carrying this filter through the
finite training geometry predicts the observed delay: reducing gamma from 64
to 8 requires \textbf{986.59 times as many updates} for the primary target.
Across five targets, measured slowdowns range from 3.34 to 986.59; all 19
available executed acquisition times lie inside the predicted intervals, with
one additional run censored. The mechanism is exact; the delay depends on
the target, geometry, and optimizer.

\section{The kernel gives gamma a direct role in optimization}
Freeze the common slope $\gamma>0$ and centers $c_j$ in the model
\begin{equation}
 f_{\gamma,\theta}(x)=b+\sum_{j=1}^W w_j\tanh\!\bigl(\gamma(x-c_j)\bigr),
 \qquad \theta=(b,w_1,\ldots,w_W).
\end{equation}
Train only the raw readout coefficients from zero on half empirical mean
squared error at $m$ training inputs $x_i$. Compare slopes with the same
centers, inputs, target, coefficient coordinates, and specified step-size rule.
For the feature vector
$\phi_\gamma(x)=(1,\tanh(\gamma(x-c_1)),\ldots,\tanh(\gamma(x-c_W)))^T$,
define
\begin{equation}
 k_\gamma(x,x')=\phi_\gamma(x)^T\phi_\gamma(x'),\qquad
 (K_\gamma)_{i\ell}=k_\gamma(x_i,x_\ell)/m.
\end{equation}
Writing $r_n$ for target values minus predictions after $n$ updates, ordinary
readout gradient descent (GD) satisfies
\begin{equation}\label{eq:update}
 \boxed{r_{n+1}=(I-\eta_\gamma K_\gamma)r_n.}
\end{equation}
Thus the kernel describes which errors a coefficient update corrects together.
Similar feature responses make two predictions move together; learning their
difference can be slow. Gamma's optimization effect follows from how it changes
this update matrix.

\section{An explicit gamma filter}
The fixed reference dictionary consists of sharp steps at the same centers:
\begin{equation}\label{eq:reference}
\begin{split}
 k_{\mathrm{step}}(x,x')&=1+\sum_{j=1}^W\sign(x-c_j)\sign(x'-c_j),\\
 \rho_\gamma(t)&=\frac{\gamma}{2}\sech^2(\gamma t).
\end{split}
\end{equation}
The density $\rho_\gamma$ averages a sharp step into a tanh transition with
width proportional to $1/\gamma$. The following theorem identifies the
resulting filter exactly.

\clearpage
\begin{theorem}[Gamma-controlled smoothing]\label{thm:mechanism}
For any fixed centers and common slope $\gamma>0$, the tanh readout kernel is
the step kernel averaged over both inputs:
\begin{equation}\label{eq:smoothing}
 k_\gamma(x,x')=\iint_{\mathbb R^2}\rho_\gamma(x-t)\rho_\gamma(x'-s)
 k_{\mathrm{step}}(t,s)\,dt\,ds.
\end{equation}
In either input, averaging multiplies a frequency-$\omega$ component by
\begin{equation}\label{eq:multiplier}
 \boxed{M_\gamma(\omega)=\frac{z}{\sinh z},\qquad
 z=\frac{\pi|\omega|}{2\gamma},\qquad M_\gamma(0)=1.}
\end{equation}
A common-slope cap $0<\gamma\le\Gamma$ implies
$M_\gamma(\omega)\le M_\Gamma(\omega)$. Frequencies well below gamma
are nearly preserved; those well above gamma have exponentially small
multipliers, $M_\gamma(\omega)\sim2ze^{-z}$ as $z\to\infty$.
\end{theorem}

\textbf{The mechanism is selective attenuation of spatial scales.}
Reducing gamma from 64 to 8 changes the multiplier at $\omega=2\pi$ from
0.9960 to 0.7851, but at $\omega=10\pi$ from 0.9074 to 0.02584. The finer
component's multiplier becomes about 35.1 times smaller. Gamma changes the
relative strength of corrections even when the step compensates for overall
kernel magnitude. Appendix~\ref{app:proof} gives the full proof, including
the transform calculation. No periodic sample or center assumption is needed.

Finite sampling and center geometry determine how the filtered components
combine into learning directions. Fourier waves need not be eigenvectors of
the sampled kernel, so $M_\gamma^2$ is generally neither a finite-kernel
eigenvalue nor an inverse learning time. We retain these couplings when
computing the prediction. The cap guarantees feature attenuation; the target's
interaction with the resulting kernel determines the delay.

\subsection*{A small example with perfect capacity}
Take one centered neuron, samples $-s,+s$ with $s>0$, and labels $-1,+1$.
Every positive gamma fits exactly with $b=0$ and $w=1/\tanh(\gamma s)$.
Yet at step $\eta=0.5$, reaching 1\% residual takes \textbf{925 updates}
when $\gamma s=0.1$ and \textbf{14 updates} when $\gamma s=1$.
The largest kernel eigenvalue is one in both cases. The target's contrast
direction has eigenvalue $\tanh^2(\gamma s)$, explaining the different
rates. Appendix~\ref{app:time} gives its exact acquisition count and the
general GD corollary.

\subsection*{What the theorem contributes to the quantitative explanation}
The theorem isolates gamma in a known frequency filter. Applying that filter
to the fixed finite geometry and target gives a learning curve through
standard least-squares GD. A controlled finite expansion makes the calculation
practical and supplies acquisition intervals. No rate is fitted to an
optimization trajectory. The following figures test the calculation
retrospectively against archived training runs.

\clearpage
\section{The filter retains the measured timing accuracy}
The primary problem has $W=559$ hidden neurons, $m=8{,}193$ equally spaced
training inputs on $[-1,1]$, equispaced centers with a boundary halo, and target
\begin{equation}\label{eq:target}
 f^\star(x)=\sin(2\pi x)+\tfrac12\sin(6\pi x)+\tfrac14\sin(10\pi x).
\end{equation}
Readouts start at zero. Each saved step is approximately $\eta_\gamma=0.5/L_\gamma$,
where $L_\gamma$ is the largest eigenvalue of $K_\gamma$. All four models
actually attain 1\% relative residual, equivalent to relative squared loss
$10^{-4}$. Their different acquisition times therefore demonstrate an
optimization effect at an attainable accuracy.

\begin{figure}[h!]
 \centering\includegraphics[width=\linewidth]{paper_three_panel.pdf}
 \caption{\textbf{Observed scale gap, predicted delay, and a slope intervention.}
 All panels use the sine-mixture target. \textbf{A:} After 20,000 joint Adam
 updates, faint points show the median and maximum absolute hidden slope
 for each of five seeds; solid curves show their medians across seeds.
 The dashed $\gamma=N/8$ construction reference is sufficient in that
 construction, not an asserted necessary threshold.
 \textbf{B:} Filter-derived intervals match executed zero-initialized raw
 readout GD hits at 1\% residual for 559 fixed-center features. The inset
 divides each interval by its observed hit. These are deterministic bounds.
 \textbf{C:} After 200,000 raw readout updates, fixed gamma 4 leaves residual
 near 0.42, while gamma $N/8$ gives 0.0034--0.0016 across widths 153--1089.
 Panels B/C use curvature-normalized steps; A motivates the scale question
 from joint training. The theorem concerns common-slope frozen features.}
 \label{fig:main}
\end{figure}

The measured first-hit time changes from 15,798,313 updates at gamma 8 to
16,013 at gamma 64. The predicted delay ratio is \textbf{985.70--987.49},
enclosing the executed \textbf{986.59}. At gamma 8, the interval's total
width is 0.181\% of the observed time. These primary endpoints have independent
interval-arithmetic certificates.

A scalar control that rescales only the gamma-64 kernel to match each case's
largest eigenvalue predicts 16,013 updates at every gamma under the saved
step rule. The full frequency-dependent filter accounts for the large delay
this magnitude-only control misses. The numerical intervals are tight for
the specified geometry and target; the theorem does not assert a universal
slowdown factor or monotone learning time for every target.

\clearpage
\section{The delay depends on the target}
Four control targets use the same geometry and four slopes: the single sine
$\sqrt2\sin(2\pi x)$, Runge function $1/(1+25x^2)$, quadratic $\sqrt5x^2$,
and exponential of sine $\exp(\sin(3\pi x))$. Every target reaches 1\%
at gamma 8 and gamma 64, allowing the measured comparison below.

\begin{figure}[h!]
 \centering\includegraphics[width=\linewidth]{target_slowdown.pdf}
 \caption{\textbf{The same gamma intervention gives target-dependent delays.}
 Each dot is the executed first-hit count at gamma 8 divided by that at
 gamma 64. Whiskers show the ratio interval from the predicted acquisition
 bounds and are narrower than the dots at this scale. Samples, centers,
 zero initialization, 1\% residual tolerance, and curvature-normalized step
 rule are shared. Only the sine-mixture endpoints have independent
 interval-arithmetic certification; the controls have numerical sensitivity
 checks. Figure~\ref{fig:all} includes every intermediate gamma and magnifies
 interval tightness.}
 \label{fig:slowdown}
\end{figure}

The same filter produces very different delays because each target places
different energy in the finite kernel's learning directions. These observations
support a quantitative target-dependent explanation. The quadratic still
needs 58.17 times as many updates at gamma 8: low polynomial degree alone
does not imply fast acquisition in this tanh dictionary. Finite geometry
and target alignment remain part of the calculation.

\subsection*{What scaling gamma with width recovers}
Figure~\ref{fig:main}C tests an intervention suggested by the filter: increase
gamma as center spacing shrinks to preserve access to finer spatial scales.
Gamma proportional to inverse spacing keeps the multiplier unchanged at
frequencies measured relative to that spacing. In this construction, that
means gamma proportional to width; Appendix~\ref{app:width} states this
precisely.

Panel C holds the physical target fixed. It demonstrates a finite-budget
benefit, and its residual near 0.42 is not an established irreducible floor.
The filter identity does not by itself guarantee constant training times
across widths. All figures concern training behavior, with no held-out
generalization claim.

\clearpage
\appendix
\section{Complete proof of the gamma mechanism}\label{app:proof}
The density in~\eqref{eq:reference} has unit mass and cumulative integral
$\int_{-\infty}^a\rho_\gamma(t)\,dt=[1+\tanh(\gamma a)]/2$.
Splitting at a step's center gives
\begin{equation}
 \int_{\mathbb R}\rho_\gamma(x-t)\sign(t-c)\,dt
 =\tanh\!\bigl(\gamma(x-c)\bigr).
\end{equation}
Insert the finite sum defining $k_{\mathrm{step}}$ into the double integral
in~\eqref{eq:smoothing}. Each pair of steps becomes a pair of tanh features,
and the constant stays one. Bounded features and an integrable density
justify all exchanges; the steps themselves need not be integrable.

Use Fourier convention
$\widehat\rho(\omega)=\int_{\mathbb R}\rho(t)e^{-i\omega t}\,dt$.
For a real number $a>0$, define
\begin{equation}
 I(a)=\int_{\mathbb R}\frac{e^{-iau}}{\cosh^2u}\,du,
 \qquad F(\zeta)=\frac{e^{-ia\zeta}}{\cosh^2\zeta}.
\end{equation}
Integrate $F$ counterclockwise around the rectangle with vertices
$-R,R,R+i\pi,-R+i\pi$, then let $R\to\infty$. On either vertical side,
$|e^{-ia\zeta}|\le e^{a\pi}$ and
$|\cosh(\pm R+iv)|^2=\sinh^2R+\cos^2v\ge\sinh^2R$
for $0\le v\le\pi$, so those integrals vanish. Since
$F(u+i\pi)=e^{a\pi}F(u)$ and the upper side runs right to left, the
horizontal integrals tend to $(1-e^{a\pi})I(a)$.

The only enclosed pole is the double pole $\zeta_0=i\pi/2$. Set
$w=\zeta-\zeta_0$. Using $\cosh(\zeta_0+w)=i\sinh w$,
\begin{equation}
 F(\zeta_0+w)=e^{a\pi/2}\left(-\frac1{w^2}+\frac{ia}{w}+O(1)\right),
 \qquad \Res_{\zeta=\zeta_0}F=ia e^{a\pi/2}.
\end{equation}
The residue theorem yields
\begin{equation}
 (1-e^{a\pi})I(a)=-2\pi a e^{a\pi/2},
 \qquad I(a)=\frac{\pi a}{\sinh(\pi a/2)}.
\end{equation}
The sine part of the real-line integrand is odd, so $I$ is real and even;
the formula extends to negative $a$. At zero,
$I(0)=\int_{\mathbb R}\sech^2u\,du=2$, the continuous limit.
Substituting $u=\gamma t$ gives the normalized multiplier
\begin{equation}
 \widehat\rho_\gamma(\omega)=\tfrac12 I(\omega/\gamma)
 =\frac{\pi|\omega|/(2\gamma)}{\sinh(\pi|\omega|/(2\gamma))}
 =M_\gamma(\omega),\qquad M_\gamma(0)=1.
\end{equation}
In particular,
\begin{equation}
 \int_{\mathbb R}\rho_\gamma(x-t)e^{i\omega t}\,dt
 =M_\gamma(\omega)e^{i\omega x}.
\end{equation}
This proves the frequency interpretation in either kernel input without
requiring an ordinary integrable Fourier transform of the steps.
Expanding $\sinh z$ gives $M_\gamma=1-z^2/6+O(z^4)$ at small $z$;
$\sinh z\sim e^z/2$ gives $M_\gamma\sim2ze^{-z}$ at large $z$.

Finally, the derivative of $z/\sinh z$ has numerator
$\sinh z-z\cosh z<0$ for $z>0$: this numerator vanishes at zero and has
derivative $-z\sinh z<0$. Reducing gamma increases $z$, proving the cap
inequality. This completes Theorem~\ref{thm:mechanism}.

\clearpage
\section{Standard GD corollary and acquisition intervals}\label{app:time}
Let $y_i=f^\star(x_i)$ be the nonzero target vector, and let $\Phi_\gamma$
have rows $\phi_\gamma(x_i)^T$. Thus $K_\gamma=\Phi_\gamma\Phi_\gamma^T/m$.
From zero initialization, with $0<\eta_\gamma\le1/L_\gamma$, define
\begin{equation}
 E_\gamma(n)=\frac{\|(I-\eta_\gamma K_\gamma)^ny\|}{\|y\|},\qquad
 n_\epsilon(\gamma)=\min\{n\in\mathbb N_0:E_\gamma(n)\le\epsilon\}.
\end{equation}
Here $0<\epsilon<1$; an empty crossing set has time $+\infty$. This is
a kernel-and-target calculation requiring no executed GD trajectory.

\begin{corollary}[Target-dependent acquisition]
Let $u_i(\gamma)$ be orthonormal eigenvectors for the positive eigenvalues
$\lambda_i(\gamma)$ of $K_\gamma$, with
$p_i(\gamma)=|u_i(\gamma)^Ty|^2/\|y\|^2$. Let
$E_{\mathrm{floor}}(\gamma)$ be the minimum relative residual over all
readouts in this same tanh model. Then
\begin{equation}\label{eq:law}
 E_\gamma(n)^2=E_{\mathrm{floor}}(\gamma)^2+
 \sum_i p_i(\gamma)(1-\eta_\gamma\lambda_i(\gamma))^{2n}.
\end{equation}
\end{corollary}
\begin{proof}
Expand $y$ in an orthonormal eigenbasis of the symmetric kernel. Nullspace
components remain unchanged and give the least-squares floor. Every positive
mode decays by the corresponding factor in~\eqref{eq:law}. Orthogonality
adds their squared residual norms.
\end{proof}
The weights are target energy fractions, or squared cosine alignments with
unit eigenvectors. Gamma changes the rates and these weights through the
explicit filter. Eigenvalues $\lambda_i$ are distinct from the dimensionless
bandwidth $\beta$ in Appendix~\ref{app:width}.

In the two-point example, the target occupies the contrast mode with
eigenvalue $\tanh^2(\gamma s)$ and zero floor. For $0<\eta<1$,
\begin{equation}
 n_\epsilon(\gamma)=\left\lceil\frac{\log\epsilon}
 {\log(1-\eta\tanh^2(\gamma s))}\right\rceil
 \sim\frac{\log(1/\epsilon)}{\eta\gamma^2s^2}\quad(\gamma s\to0).
\end{equation}

For the full experiment, a finite harmonic expansion retains the sample/center
couplings and bounds omitted terms, including distant transitions in its
auxiliary periodic representation. It approximates the original nonperiodic
tanh model with the original readout coordinates. The detailed construction
is in the \href{\repo/docs/gamma_factorized_readout.md}{technical supplement}.

If $\widehat E_\gamma(n)$ is the predicted residual and a justified margin
satisfies $|E_\gamma(n)-\widehat E_\gamma(n)|\le d_\gamma(n)$, then
integer witnesses $a,b$ give
\begin{equation}
\begin{gathered}
 \widehat E_\gamma(a)-d_\gamma(a)>\epsilon,\qquad
 \widehat E_\gamma(b)+d_\gamma(b)\le\epsilon\\[-2pt]
 \Longrightarrow\quad a+1\le n_\epsilon(\gamma)\le b.
\end{gathered}
\end{equation}
The true residual is nonincreasing under the step condition, proving the
necessary count; the second witness proves sufficiency. The simplest margin
is $n\eta_\gamma\Delta$ for kernel error bound $\Delta$, with both update
operators contracting. The refinement retains the error's action on the
target. Dividing necessary and sufficient counts across gammas gives the
ratio interval. Numerical verification status is separate from these
exact-arithmetic statements (Appendix~\ref{app:cases}).

\clearpage
\section{Width scaling preserves relative bandwidth}\label{app:width}
For center spacing $h$, set $\beta=\gamma h$ and grid-relative frequency
$\xi=h\omega$. Fixing $\beta>0$ gives
\begin{equation}
 M_{\beta/h}(\xi/h)=
 \frac{\pi|\xi|/(2\beta)}{\sinh(\pi|\xi|/(2\beta))},
\end{equation}
independent of $h$, with value one at zero. For fixed gamma and nonzero
$\xi$, as $h\to0$,
\begin{equation}
 M_\gamma(\xi/h)^2\sim\frac{\pi^2\xi^2}{\gamma^2h^2}
 \exp\!\left(-\frac{\pi|\xi|}{\gamma h}\right).
\end{equation}
The construction uses $h=2/N$, $W=N+2\lceil\sqrt N\rceil+1$, and
$\gamma=N/8$, hence $\beta=1/4$ and $\gamma=\Theta(W)$.
This avoids a growing exponential attenuation at the dictionary's resolution.
An upper bound $\gamma=O(W)$ alone also permits fixed gamma and does not
imply the conclusion. At fixed physical frequency, fixed gamma already
gives width-independent attenuation; linear scaling makes its multiplier
approach one. Geometry and target weights still determine training time.

\section{Protocol and evidence provenance}\label{app:protocol}
Figure~\ref{fig:main}A trains hidden slopes, offsets, and readout parameters
for 20,000 Adam updates, with learning rate $10^{-3}$, epsilon $10^{-8}$,
and moments $(0.9,0.999)$. Hidden slopes and readout weights start independently
uniform on $[-\sqrt{6/(W+1)},\sqrt{6/(W+1)}]$; biases start at zero.
Seeds 0--4 are all shown. The width study uses $N=128,256,512,1024$;
the acquisition study uses $N=512$. Grids have $m=16N+1$ equally spaced
endpoint samples on $[-1,1]$. Frozen centers have spacing $2/N$ with a halo
of $\lceil\sqrt N\rceil$ additional centers at each end.

In Figure~\ref{fig:main}C, the target, inputs, centers, raw coordinates, and
200,000-update budget are shared within each width. Gamma is the intervention,
and the step follows each kernel's curvature. These are measured residuals;
certified times have not been computed at every width. Joint-training
observations in A do not establish a theorem about joint Adam dynamics.

The four primary intervals have independent 192-bit interval-arithmetic
endpoint certificates for nominal-real tanh on archived binary inputs and
saved steps. They do not certify every floating-point iterate or plotted
curve. The sixteen control intervals have double-precision reference and
arithmetic-sensitivity checks, including a disclosed heuristic allowance for
transcendental evaluations. No separate interval certificates were computed
for those targets. A zero-width numerical interval alone does not certify
an exact first hit.

\textbf{Review materials.}
The \href{\repo/docs/gamma_optimization_paper_note.md}{companion Markdown note}
and \href{\repo/docs/gamma_factorized_readout.md}{technical supplement}
provide the narrative and full numerical construction.
The \href{\evidence/REPORT.md}{evaluation report} includes control targets,
the scalar-rescaling control, analytic versus combined margins, and sensitivity
checks. The \href{\evidence/summary.json}{numerical summary},
\href{\evidence/interval_audit.json}{interval audit}, and
\href{\evidence/review_figure_data.json}{figure data} preserve exact values
and verification status. Figures reuse archived results; no new training was
needed. The \href{\evidence/paper_validation.json}{validation record}
separates prior implementation tests from this document's checks.

\clearpage
\section{Complete five-target acquisition results}\label{app:cases}
All cases share the 559-neuron geometry, 8,193 samples, zero raw readout,
half empirical mean squared error, and each gamma's saved step, approximately
$0.5/L_\gamma$. Intervals combine the analytic and target-action margins.
Harmonic resolution is selected from calculated bounds, without fitting
endpoints to executed hits.

\begin{figure}[h!]
 \centering\includegraphics[width=\linewidth]{target_acquisition.pdf}
 \caption{\textbf{All twenty predictions, with nineteen executed crossings.}
 \textbf{A--E:} Colored curves and whiskers give predicted first-hit intervals;
 black dots give executed hits at 1\% residual. Vertical ranges differ;
 connecting lines guide the eye. The open triangle marks the quadratic run
 stopped at 200,000 updates and is not an observed crossing.
 \textbf{F:} Each interval endpoint is divided by its executed hit, shifted
 by one, and multiplied by 100. The symmetric-log axis expands small
 deviations. All 19 intervals contain zero; the censored case is excluded.
 The largest total relative width is 0.181\%, for the sine mixture at gamma 8.
 Only A's intervals have independent interval-arithmetic certificates.}
 \label{fig:all}
\end{figure}

The censored quadratic run at gamma 12 has a predicted crossing interval
269,286--269,299 and a separate original-kernel spectral forecast of 269,292.
Both remain forecasts. The other nineteen cases compare predictions against
actual GD crossings. These checks assess the target-dependent timing
calculation on archived training problems; they do not establish a universal
learning-time ordering over targets or dictionaries.
\end{document}
